% ---- Acronyms used throughout the dissertation -----------------------------
\newacronym{gfet}{GFET}{Graphene Field-Effect Transistor}
\newacronym{fet}{FET}{Field-Effect Transistor}
\newacronym{ssdna}{ssDNA}{single-stranded DNA}
\newacronym{msa}{MSA}{Multiple Sequence Alignment}
\newacronym{ppi}{PPI}{Percentage of Pairwise Identity}
\newacronym{ncbi}{NCBI}{National Center for Biotechnology Information}
\newacronym{tm}{$T_m$}{melting temperature}
\newacronym{mfe}{MFE}{Minimum Free Energy}
\newacronym{amr}{AMR}{Antimicrobial Resistance}
\newacronym{poc}{PoC}{Point-of-Care}
\newacronym{plddt}{pLDDT}{predicted Local Distance Difference Test}
\newacronym{pae}{PAE}{Predicted Aligned Error}

\chapter{Introduction}

\section{Context and motivation}

The timely identification of pathogenic microorganisms is a cornerstone of clinical
diagnostics, public-health surveillance and the containment of antimicrobial resistance
(\acrshort{amr}). Conventional diagnostic methods, such as microbiological culture and
polymerase chain reaction, are sensitive and well established, but they are typically slow,
laboratory-bound and require trained personnel and specialised equipment. There is therefore a
strong demand for \acrfull{poc} devices that are fast, portable, inexpensive and able to
detect a pathogen directly from a biological sample with minimal preparation.

Biosensors based on the \acrfull{gfet} have emerged as one of the most attractive platforms to
meet this demand. Graphene, a single layer of carbon atoms arranged in a honeycomb lattice,
combines exceptional electrical conductivity with an entirely surface-exposed conduction
channel, so that any charged molecule binding at its surface produces a measurable change in
the transistor current \citep{novoselov2004,schedin2007}. When the graphene channel is
functionalised with a \acrfull{ssdna} capture probe, the hybridisation of a complementary
target strand can be transduced into an electrical signal in a label-free manner and with very
high sensitivity, down to the single-molecule regime in optimised devices
\citep{hwang2020,beraud2021}. This makes the \acrshort{gfet} an ideal transducer for
nucleic-acid-based detection of pathogens.

The analytical performance of a \acrshort{gfet} biosensor is, however, only as good as the
capture probe immobilised on its surface. A useful probe must simultaneously satisfy several,
partly competing, requirements. It must be complementary to a region of the target genome that
is \emph{conserved} across the clinical strains of the pathogen, otherwise strain-to-strain
variation would cause false negatives. It must hybridise with its target strongly and
specifically at the temperature of the assay, which constrains its length, melting temperature
and GC content. Finally, it must not fold back on itself into stable secondary structures
(hairpins or self-dimers), because a probe that is folded is not available to capture the
target. Balancing these criteria for a single probe is already non-trivial; doing so
systematically for many genes and organisms, while ranking the surviving candidates in a
reproducible and quantitative way, is beyond what manual design can offer.

\section{Aims and objectives}

The goal of this dissertation is to design, \emph{in silico}, high-quality \acrshort{ssdna}
capture probes for a panel of clinically relevant bacterial pathogens, to be used in the
fabrication of \acrshort{gfet} biosensors. To that end, a single, reproducible computational
pipeline was developed that takes a target gene as input and returns a small set of ranked,
biophysically validated probe candidates. The specific objectives were:

\begin{itemize}
  \item to automate the retrieval and curation of gene sequences from the \acrshort{ncbi}
        databases, and their alignment, so that conserved regions can be identified without
        manual intervention;
  \item to screen candidate probes with well-founded biophysical criteria --- sequence
        conservation, melting temperature, GC content, and hairpin, homodimer and secondary
        structure stability --- and to combine these into transparent quality scores;
  \item to make the design criteria \emph{organism-aware}, so that targets with very different
        genome composition (for example AT-rich or GC-rich pathogens) are evaluated fairly;
  \item to validate the most promising candidates by three-dimensional structure prediction;
        and
  \item to integrate and compare the computationally designed probes with an existing
        experimental probe panel, and to produce curated datasets for downstream use.
\end{itemize}

\section{Target panel}

Six marker genes from six clinically important bacterial species were used as the working
panel throughout this work (Table~\ref{tab:targets}). They were selected as established
species-specific markers and span a wide range of genome composition, from AT-rich
(\emph{Staphylococcus aureus}, \emph{Haemophilus influenzae}) to GC-rich (\emph{Pseudomonas
aeruginosa}), which makes them a demanding test of a general-purpose design method.

\begin{table}[H]
\begin{center}
\begin{tabular}{l l c l}
\hline
\textbf{Gene} & \textbf{Organism} & \textbf{Group} & \textbf{Function} \\
\hline
\emph{nuc}  & \emph{Staphylococcus aureus}    & A & thermonuclease \\
\emph{rmpM} & \emph{Neisseria meningitidis}   & A & class-4 outer-membrane protein \\
\emph{lytA} & \emph{Streptococcus pneumoniae} & B & autolysin \\
\emph{oprL} & \emph{Pseudomonas aeruginosa}   & B & peptidoglycan-associated lipoprotein \\
\emph{algD} & \emph{Pseudomonas aeruginosa}   & B & GDP-mannose 6-dehydrogenase \\
\emph{frdB} & \emph{Haemophilus influenzae}   & B & fumarate reductase Fe--S subunit \\
\hline
\end{tabular}
\end{center}
\caption{Target genes and their host organisms used as the working panel.}
\label{tab:targets}
\end{table}

\section{Structure of the dissertation}

The remainder of this dissertation is organised as follows. Chapter~2 reviews the state of the
art in graphene biosensors and in the computational design of oligonucleotide probes.
Chapter~3 describes the methods, detailing every stage of the pipeline and the rationale for
each design decision and threshold. Chapter~4 presents and discusses the results obtained for
the target panel, including the screening funnel, the three-dimensional validation, the
diversity and rarefaction analyses, and the comparison with the reference probe panel.
Chapter~5 draws the main conclusions and outlines directions for future work.
